% Lima kualitas sistem yang saling tarik-menarik. Setiap sisi segi lima
% menandai satu pasangan yang sering bertentangan, dengan contoh singkatnya.
\begin{tikzpicture}[
  font=\scriptsize,
  kualitas/.style={draw, rounded corners, align=center, minimum height=0.85cm,
                   text width=2.1cm, fill=blue!8, font=\scriptsize\bfseries},
  tarik/.style={{Stealth[length=1.8mm]}-{Stealth[length=1.8mm]}, thick, red!60!black,
                dashed},
  catatan/.style={font=\tiny, gray!40!black, align=center}
]
  \node[kualitas] (sedia) at (90:3.1) {Ketersediaan};
  \node[kualitas] (biaya) at (18:3.1) {Biaya};
  \node[kualitas] (skala) at (-54:3.1) {Skalabilitas};
  \node[kualitas] (aman) at (-126:3.1) {Keamanan};
  \node[kualitas] (ubah) at (162:3.1) {Kecepatan\\perubahan};

  \node[draw, rounded corners, align=center, fill=praditagreen!15, text width=2.9cm,
        minimum height=1.1cm, font=\scriptsize] at (0,0)
    {\textbf{Keputusan arsitektur}\\[2pt]
     {\tiny\mdseries memilih titik seimbang\\dan menuliskan alasannya}};

  \draw[tarik] (sedia) -- node[catatan, above right=-1pt] {cadangan di dua lokasi\\menambah biaya} (biaya);
  \draw[tarik] (biaya) -- node[catatan, right=2pt] {kapasitas cadangan\\tetap dibayar} (skala);
  \draw[tarik] (skala) -- node[catatan, below=0.55cm] {lebih banyak komponen,\\lebih banyak celah} (aman);
  \draw[tarik] (aman) -- node[catatan, left=2pt] {pemeriksaan tambahan\\memperlambat rilis} (ubah);
  \draw[tarik] (ubah) -- node[catatan, above left=-1pt] {perubahan sering\\membuka peluang gangguan} (sedia);
\end{tikzpicture}
